\begin{tikzpicture}[
  font=\small,>={Stealth[length=2mm]},
  input/.style={draw=blue!60!black,fill=blue!7,rounded corners=2pt,
    minimum width=29mm,minimum height=10mm,align=center},
  path/.style={draw=orange!75!black,fill=orange!10,rounded corners=2pt,
    minimum width=31mm,minimum height=11mm,align=center},
  meter/.style={draw=teal!70!black,fill=teal!8,rounded corners=2pt,
    minimum width=35mm,minimum height=11mm,align=center},
  flow/.style={->,thick,draw=gray!80}
]
\node[input] (bus) at (0,2.25) {公共母线\\送端功率};
\node[input] (tank) at (0,0.75) {合格氢库存\\可提取量};
\node[input] (work) at (0,-0.75) {完成任务\\有效服务量};
\node[input] (demand) at (0,-2.25) {海洋用能需求};
\node[path] (cable) at (4.0,2.25) {海缆容量、可用率\\损耗与陆侧接纳};
\node[path] (logistics) at (4.0,0.75) {管输/船运容量\\窗口与产品损失};
\node[path] (fiber) at (4.0,-0.75) {光纤容量、时延\\与SLA审核};
\node[path] (marine) at (4.0,-2.25) {刚柔性负荷\\与未供量};
\node[meter] (mpE) at (8.3,2.25) {海缆陆上受端\\结算电量};
\node[meter] (mpH) at (8.3,0.75) {氢气接收端\\结算质量};
\node[meter] (mpC) at (8.3,-0.75) {算力服务端\\可结算服务量};
\node[meter] (mpM) at (8.3,-2.25) {海洋设备入口\\实际服务量};
\draw[flow] (bus) -- (cable); \draw[flow] (cable) -- (mpE);
\draw[flow] (tank) -- (logistics); \draw[flow] (logistics) -- (mpH);
\draw[flow] (work) -- (fiber); \draw[flow] (fiber) -- (mpC);
\draw[flow] (demand) -- (marine); \draw[flow] (marine) -- (mpM);
\end{tikzpicture}
